\documentclass[11pt]{article}
\usepackage{series}
\newcommand{\frozenresult}[2]{\href{https://github.com/PeterNero/mtt-results-repro/blob/f141a20ea23c5c3ff19cc2161c0e226e29ade8a7/release/results/#1/artifact.json}{#2}}
\usepackage{array,tabularx}

\title{Closure-Strain Geometry:\\
Local Normal Forms and Conditional Matter Encodings}
\author{Peter Nero}
\date{September 2026\\Version 8}

\begin{document}
\maketitle

\begin{abstract}
We extract the exact geometry of closure strain from earlier Standard-Model
interpretations.  For an invertible rank-three comparison field, polar or
Iwasawa reduction removes three orientation directions and leaves six strain
directions.  With a selected orthonormal flag these split orthogonally into
scalar, traceless-diagonal, and shear sectors of dimensions $1+2+3$.  We give
explicit projectors and a norm identity, explain the relation and distinction
between symmetric shear and the Heisenberg nil algebra, and state precisely
what a closure Hessian proves.  In particular, Hessian positivity does not
select a unique Higgs, three families, Standard Model charges, confinement,
mixing, or CP violation.  Those identifications require representation,
connection, action, and source theorems.  Current q79 and finite-algebra
calculations are incorporated at their declared profile-equivalence tier, and
the q79-side finite bridge is sharpened: one universal flat differential line
commutes with the $1+2+3$ lane projectors and lifts the normalized Reynolds
Hessian exactly. This closes connection/holonomy and Hessian naturality at
finite-symbol tier. The same-source continuum intertwiner from the local
strain normal form to the nonzero-Chern physical HYM carrier remains isolated
and open. The selected internal TT co-shape support is explained separately:
its helicity-two carrier and normalized internal eigenvalue are exact, while
the physical stress-energy map and SI normalization are not supplied by them.
The curved connection compiler and nonlinear spectral-strain shorting now
replace the excluded linear root-space route, including the later
support-stratified boundary law. Finite higher transfer is distinguished from
both an analytic fixed-point theorem and a physical strain action.
\end{abstract}

\section*{Revision note for version 8}
\addcontentsline{toc}{section}{Revision note for version 8}
\begin{description}
\item[Supersedes.] Version~7; the preceding edition's revision note is retained
below.
\item[Reason.] The frozen TT support result appeared only in the evidence
list, leaving its relation to strain, helicity, and physical normalization
unexplained.
\item[Resolution.] Version~8 imports the scoped support argument and explains
the distinction between internal support, physical stress-energy response,
cost pullback, and reducing operator transport. It also explains the curved
connection compiler, nonlinear strain shadow, support-stratified boundary,
finite transfer hierarchy, and real/complex and polar-metric qualifications.
\item[Retained.] The local normal form, finite shared-line/Hessian square,
established internal TT support, and finite/profile SM conclusions retain
their original tiers.
\item[Open boundary.] Physical nonflat comparison, stress-energy response and
SI normalization remain open. Established support and finite/profile results
are retained.
\end{description}

\section*{Revision note for this edition}
\addcontentsline{toc}{section}{Revision note for this edition}
\begin{description}
\item[Supersedes.] \emph{Closure-Strain Geometry: Local Normal Forms and
Conditional Matter Encodings}, version~6.
\item[Reason.] Version~6 correctly left the physical local-to-q79 intertwiner
open, but it predates the exact q79 universal-line and finite-Hessian square
and therefore understates what is already closed on the target side.
\item[Resolution.] Version~7 retains the local $1+2+3$ theorem and adds the
flat differential-line pullback, finite Reynolds projector, exact Hessian
spectrum, and the precise curved-HYM nonpromotion guard.
\item[Retained result.] The local strain normal form and the executed embedded
renormalized-SM profile branch both survive, as distinct results.
\item[Remaining boundary.] Their physical identification still requires the
same-source metric, connection, and continuum-HYM intertwiner; strict no-knob
branch and value selection remain stronger upgrades.
\end{description}

\tableofcontents

\section{How to Read Closure Strain}

Closure strain is best understood first as a local stiffness problem.  The
comparison field $Q$ measures how two rank-three frames differ.  Polar
decomposition removes the orientation part, and $S=\log U$ records the
remaining deformation.  A closure functional $\mathcal J$ then assigns a cost
to perturbing that aligned state.

The paper has three distinct achievements:
\begin{enumerate}
\item the local strain space has an exact flag-selected $1+2+3$ orthogonal
normal form;
\item the selected q79 target has an exact finite rank-$1+2+3$ carrier,
shared line, and finite Hessian square; and
\item the finite Standard Model construction realizes matter and scalar
profiles at its declared profile tier.
\end{enumerate}
These results are mutually compatible, but they become one derivation only
after a same-source bundle, connection, metric, and operator intertwiner is
constructed.

\subsection{Stiffness is not yet mass}

For a finite model, write the quadratic action schematically as
\[
 \frac12\dot s^{T}K\dot s-\frac12s^{T}Hs.
\]
The closure Hessian $H$ measures restoring cost.  The physical squared
frequencies or masses solve the generalized eigenvalue problem
\[
 Hv=m^2Kv.
\]
Rescaling a field changes the matrix entries of both $H$ and $K$ while leaving
the physical generalized eigenvalues unchanged.  Consequently a positive
entry of $H$, or even a positive eigenvalue computed with an arbitrary
coordinate norm, is not yet a mass prediction.

The same distinction applies to particle names.  A one-dimensional scalar
strain block is a geometric representation.  To call it the Higgs requires a
complex weak doublet representation, a selected kinetic term and potential,
symmetry breaking, and the renormalized pole map.  The local dimension count
only makes such an intertwiner possible.

\subsection{Argument map}

Sections~2--6 establish the status, comparison, projector, Iwasawa, and
Hessian geometry.  Sections~7--10 state the independent Higgs, matter, family,
Yukawa, mixing, and CP gates.  Section~11 identifies what is already closed on the q79
target side and isolates the continuum source map still needed to join the
local geometry to that target.

\section{Purpose and status}

The paper proves a local normal form and organizes possible matter encodings.
It does not derive the Standard Model from the normal form.  We distinguish:
\begin{description}
\item[Local theorem.] Linear algebra and differential geometry following from
the stated rank-three comparison field and selected flag.
\item[Selected realization.] The q79 trace carrier and the finite
$\mathbb C\oplus\mathbb H\oplus M_3(\mathbb C)$ branch supplied by independent
packets.
\item[Profile equivalence.] Reproduction of the embedded renormalized Standard
Model using the declared one-shared-physical-primitive/profile standard.
\item[Strict selection.] Derivation of the observed branch and numerical data
without target-profile selection; this remains a stronger program.
\end{description}

\section{Comparison field and strain}

Let $TP,TI\to B$ be oriented Euclidean rank-three bundles and let
\[
 Q=Q_{\rm WW}\in\Gamma(\operatorname{Hom}(TP,TI))
\]
be invertible and orientation preserving on an admissible domain.  In local
orthonormal frames,
\[
 Q=RU,
 \qquad R\in SO(3),
 \qquad U=(Q^TQ)^{1/2}>0.
\]

\begin{definition}[Logarithmic closure strain]
The logarithmic strain is
\[
 S=\log U=\frac12\log(Q^TQ)
 \in\operatorname{Sym}(3,\mathbb R).
\]
It is dimensionless.  Any physical dimension assigned to a fluctuation of $S$
must come from a separately normalized action and field coordinate.
\end{definition}

The nine local comparison components split as
\[
 \operatorname{Mat}(3,\mathbb R)
 =\mathfrak{so}(3)\oplus\operatorname{Sym}(3,\mathbb R),
 \qquad 9=3+6.
\]
This is a local field-component decomposition, not a product decomposition of
a manifold.
Its orthogonality is the skew/symmetric split at the identity, not a claim
that polar variations $\Omega U$ and $\delta U$ are orthogonal in the ambient
Frobenius metric at an arbitrary $U$. The logarithmic strain has its own
declared metric; transporting a cost to these coordinates requires its
pullback, including the background-dependent derivative.

\section{Exact \texorpdfstring{$1+2+3$}{1+2+3} normal form}

Choose an orthonormal flag.  For $S\in\operatorname{Sym}(3,\mathbb R)$ write
\[
 \operatorname{Diag}(S)=\text{the diagonal matrix with the diagonal of }S
\]
and define
\[
 P_{\rm sc}S=\frac{\operatorname{tr}S}{3}I_3,
 \qquad
 P_{\rm sh}S=\operatorname{Diag}(S)-P_{\rm sc}S,
 \qquad
 P_{\rm nil}S=S-\operatorname{Diag}(S).
\]

\begin{theorem}[Closure-strain projector theorem]
The three maps are self-adjoint orthogonal projectors for the Frobenius inner
product.  Their images have dimensions $1,2,3$, and
\[
 S=P_{\rm sc}S+P_{\rm sh}S+P_{\rm nil}S,
 \qquad
 \|S\|_F^2
 =\|P_{\rm sc}S\|_F^2
 +\|P_{\rm sh}S\|_F^2
 +\|P_{\rm nil}S\|_F^2.
\]
\end{theorem}

\begin{proof}
The images are, respectively, scalar diagonal matrices, traceless diagonal
matrices, and symmetric zero-diagonal matrices.  Their pairwise Frobenius
inner products vanish.  Each displayed map is the identity on its image and
zero on the other two images, proving idempotence and self-adjointness.  Their
dimensions are one, two, and three, which sum to
$\dim\operatorname{Sym}(3)=6$.
\end{proof}

\begin{remark}[Dependence on the flag]
The trace line is canonical, but the split of the traceless five-dimensional
sector into dimensions two and three depends on the selected flag.  Under the
full $SO(3)$ action the invariant decomposition is $1+5$.  Therefore an MTT
claim using $1+2+3$ must identify the operator or geometry selecting the flag.
\end{remark}

\section{Iwasawa interpretation and the nil warning}

For $Q\in GL^+(3,\mathbb R)$, QR/Iwasawa decomposition gives
\[
 Q=KAN,
 \qquad K\in SO(3),
\]
where $A$ is positive diagonal and $N$ is upper unitriangular.  Their
dimensions are
\[
 \dim K=3,
 \qquad \dim A=3=1+2,
 \qquad \dim N=3.
\]
The Lie algebra of $N$ consists of strictly upper-triangular matrices and is
the three-dimensional Heisenberg nil algebra.

The symmetric off-diagonal space $\operatorname{im}P_{\rm nil}$ is also
three-dimensional, but it is not itself that nil Lie algebra: symmetric
off-diagonal matrices are not closed under commutators.  The two spaces are
linearly identified at a chosen background through the differential of the
polar/Iwasawa coordinate change.  A paper using the word ``nil'' must state
whether it means the Heisenberg group, its Lie algebra, a symmetric shear
coordinate, or a boundary/anchoring operator.

This resolves the valid core of the circle--lens--nil picture: it can describe
a phase line, a finite/projective transport lane, and a triangular anchoring
lane.  It does not prove a literal nested or product topology.

\section{Closure cost and Hessian consequences}

Let $\mathcal J$ be a $C^3$ gauge-invariant closure functional on a Sobolev
neighborhood of an aligned configuration $S_\ast$, and suppose
$D\mathcal J(S_\ast)=0$.  Its quadratic variation on a gauge-fixed slice is
\[
 \mathcal J(S_\ast+s)
 =\mathcal J(S_\ast)
 +\frac12\langle s,Hs\rangle+R_3(s),
 \qquad |R_3(s)|\le C\|s\|^3.
\]

\begin{theorem}[What a positive strain Hessian proves]
If $H$ is self-adjoint and $H\ge cI$ on the six-dimensional strain slice for
some $c>0$, then $S_\ast$ is a strict local minimum on that slice and all six
quadratic strain directions have positive restoring cost.  Positivity alone
does not select one of the six directions as a physical scalar and does not
identify any eigenvalue with a pole mass.
\end{theorem}

\begin{proof}
Taylor's theorem and the coercive quadratic bound imply strict local
minimality for sufficiently small $s$.  A positive operator can have between
one and six distinct eigenspaces and carries no representation or kinetic
normalization data.  Pole masses depend on the canonically normalized kinetic
operator as well as $H$.
\end{proof}

If $H$ commutes with all three projectors, it has block form
\[
 H=H_{\rm sc}\oplus H_{\rm sh}\oplus H_{\rm nil}.
\]
This block invariance is an additional theorem or symmetry assumption.  A
generic Hessian mixes the three selected sectors.

Even in the block-diagonal case, the spectrum answers a stability question
before it answers a particle question.  The eigenvectors must transform in
the required gauge representations, survive the constraints, and be
canonically normalized by the kinetic operator.  Only then can an eigenvalue
be transported to a pole observable.

\section{Conditional Higgs interpretation}

The trace line $\operatorname{im}P_{\rm sc}$ supplies one local scalar
coordinate.  To identify it with a Higgs degree of freedom one must additionally
prove:
\begin{enumerate}
\item the selected finite representation is an $SU(2)$ complex doublet with
the required hypercharge, rather than a real singlet strain coordinate;
\item the alignment projector selects one rank-four real doublet module and
removes or lifts all other scalar modules;
\item the kinetic term is positive and canonically normalized;
\item the potential has the required vacuum and symmetry-breaking orbit; and
\item the renormalized pole observable follows after threshold and RG
transport.
\end{enumerate}

In the selected finite-algebra computation, the unrestricted real fluctuation
space contains three rank-four scalar-doublet modules.  The q79/proto-spinor
alignment rule is represented by a rank-four projector that removes eight
unwanted real scalar directions at the declared profile tier.  This is a
substantive downstream selection result.  It is not implied by Hessian
positivity, and its identification with the local trace line still needs the
same-source intertwiner.

\section{Conditional matter and gauge encodings}

The selected finite branch uses
\[
 \mathcal A_F=\mathbb C\oplus\mathbb H\oplus M_3(\mathbb C)
\]
on an explicit three-family particle--antiparticle carrier.  At the profile
tier its finite Dirac operator satisfies the declared self-adjointness,
grading, real-structure, order-zero, and order-one checks.  The anomaly-free
circle and the usual gauge quotient are recovered on that selected spectrum.

These results are downstream finite-algebra theorems.  They do not follow from
the dimensions $1,2,3$ alone.  In particular:
\begin{itemize}
\item rank one does not by itself select hypercharge normalization;
\item rank two does not by itself select weak chirality or a quaternionic real
structure;
\item rank three does not by itself select color representations or
confinement; and
\item three carrier lanes do not by themselves prove three fermion families.
\end{itemize}

The exact representation, bimodule, grading, real structure, and anomaly
calculation must remain visible wherever Standard Model language is used.

\section{Families, quarks, and order of breakdown}

The hypothesis that quarks represent a second-order closure breakdown can be
encoded by assigning leptonic states to a first nontrivial quotient and quark
states to an iterated or noncommuting response sector.  To turn this into a
theorem, one needs a filtration
\[
 0\subset F_1\subset F_2\subset\mathcal H_{\rm q79}
\]
and a selected operator whose first response preserves $F_1$ while its second
response reaches $F_2/F_1$.  The construction must reproduce representation
content and source values without using the observed quark/lepton distinction
as the selector.  The present normal form makes such a theorem well-typed but
does not prove it.

Family multiplicity is similarly an index or monodromy problem, not a
consequence of the number of strain blocks.  The selected $\mathbb Z_3$
family/character factor and the gauge-rank flag must remain separate tensor
factors.

\section{Yukawa magnitudes, mixing, and CP}

A geometric Yukawa entry requires normalized zero modes and a trilinear
functional, schematically
\[
 Y_{ijk}=\int_{X_6}
 \langle\psi_i,\mathcal K_H(\psi_j,h_k)\rangle\,\mathrm{vol}_{X_6}.
\]
Its value depends on the selected connection, metrics, normalization, Higgs
mode, and transport scheme.  Complex phase can arise from common-circle
holonomy, but the holonomy class and orientation branch must be selected
before comparison with CKM or PMNS data.

The current repositories close charged Yukawa magnitudes, the finite matrix,
CKM profile, electroweak rows, and precision transport at the adopted embedded
renormalized-SM/one-shared-physical-primitive/profile standard.  This means the
constructed finite operator is capable of the required SM profile and passes
its frozen certificates.  It does not mean the six-dimensional strain normal
form alone derives all measured values, nor does it close strict no-knob unique
selection.

\section{The q79 connection problem}

For the selected degree-three cover,
\[
 \mathcal H_{\rm q79}
 =L_{\rm shared}\otimes
 (\mathcal O\oplus\mathcal A_0\oplus\mathcal A),
 \qquad \operatorname{rank}=1+2+3.
\]

On the q79 branch complement, the unique nontrivial map
\[
 h_{S_3}:S_3\to\mathbb Z_{64},
 \qquad h_{S_3}(\sigma)=32\,\epsilon(\sigma),
\]
pulls the universal flat weight-one line over
$B_\nabla\mathbb Z_{64}$ back to the SpinC determinant sign line for either
admissible odd root. Denote this specified pullback by
$(L_{\rm sh},\nabla_{\rm sh})$.

\begin{theorem}[Closed q79 finite target square]
The same line $L_{\rm sh}$ tensors all three q79 lanes and its scalar holonomy
commutes with their projectors. For the selected two-copy sheet/edge symbol,
\[
 P_{\rm Haar}=\frac1{6}\sum_{g\in S_3}\rho(g),
 \qquad
 H_{\rm fin}=\kappa_{\rm fin}(I-P_{\rm Haar})
\]
satisfy
\[
 \operatorname{rank}P_{\rm Haar}=2,
 \qquad
 \operatorname{spec}(H_{\rm fin}/\kappa_{\rm fin})
 =\{0^{\times2},1^{\times4}\},
 \qquad H_{\rm TT}=\kappa_{\rm fin}I_2.
\]
The lifted operator is exactly
$I_{L_{\rm sh}}\otimes H_{\rm fin}$. Thus the q79 connection/holonomy line and
finite Hessian square are closed without a dimensionless fit.
\end{theorem}

\begin{proof}
The character identity $\chi_r\circ h_{S_3}=\operatorname{sgn}$ gives the
parallel determinant-line pullback. Scalar holonomy commutes with every sheet
operator. Haar averaging is an orthogonal projector, and decomposition of the
selected representation gives the displayed ranks and spectrum.
\end{proof}

This theorem closes the target-side finite square, not the physical bridge.
The flat root-stack line cannot equal the nonzero-Chern physical HYM
connection. Likewise, it supplies no map from the local world-in-world strain
bundle.

What has been gained is a fixed codomain for the missing map.  The target
projectors, line action, and normalized finite Hessian no longer have to be
guessed while constructing the intertwiner.  The remaining task is to show
that the physical continuum geometry emits those already fixed objects, not
merely another rank-six model with a similar spectrum.

The local strain bundle has the same rank profile over a different scalar
field. For a complex-linear identification, complexify the real strain
bundle; alternatively supply a compatible real form of the q79 carrier.
The candidate map in the first convention is
\[
 \mathfrak I:
 (\operatorname{im}P_{\rm sc}\oplus
 \operatorname{im}P_{\rm sh}\oplus
 \operatorname{im}P_{\rm nil})\otimes_{\mathbb R}\mathbb C
 \longrightarrow \mathcal H_{\rm q79}.
\]

\begin{theorem}[Requirements for a linear physical identification]
Within the proposed bundle-isomorphism route, rank matching promotes to a
same-source linear identification only if
$\mathfrak I$ is a global bundle isomorphism and, on the selected domains,
\[
 \mathfrak I^*G_{\rm HYM}=G_{\rm strain},
 \qquad
 \mathfrak I\nabla^{\rm strain}=\nabla^{\rm HYM}\mathfrak I,
 \qquad
 \mathfrak I H_{\rm strain}=H_{\rm q79}\mathfrak I,
\]
with analogous identities for the retarded and overlap kernels used to emit
the finite source rows.
\end{theorem}

\begin{proof}
A linear identification must be independent of local trivialization and
must preserve the structures used to define the action and observables.
Bundle, metric, connection, and operator intertwining are therefore necessary.
They are also sufficient to transport the declared quadratic and overlap
calculations between the two descriptions.
\end{proof}

These conditions specify the strong linear route, not every possible
effective reduction. In particular, the viable spectral-strain construction
below is nonlinear and shorted. The closed finite target square survives;
an excluded linear realization must not be revived merely to match it.

\subsection{Curved geometry and the nonlinear strain shadow}
\label{sec:consumer-geometry}
The QM owner describes the physical continuum target using a covariant
projective module, retaining its connection and domain data
\cite{QMGeometry,FrozenProjective,FrozenCech}. A smooth finite matrix
projector $p(x)$ still has infinitely many possible base modes. Its explicit
\v{C}ech realization carries a correction $\Gamma$ to the Grassmann connection;
forgetting $\Gamma$ generally changes the supplied physical connection.
Naturality transports the actual connection, curvature, metric and spectral
windows together. It does not make a chosen raw Fourier window reducing.
The CLN ranks describe a local, auxiliary or filtrational carrier, not
$S^1\times\mathrm{Lens}\times\mathrm{Nil}$ as the global q79 topology.
An external flag is not a parallel subbundle inserted into an irreducible
physical HYM connection.

\label{sec:consumer-strain}
The physical relative-phase action on $\operatorname{Herm}(3)$ excludes the
old stable linear rank-six root-space identification: its $S_3$ character is
$(6,0,0)$, whereas the diagonal/edge carrier has character $(6,2,0)$.
The replacement is the nonlinear map
\[
 \Phi(X)=(d_0,d_1,d_2,|z_{12}|^2,|z_{20}|^2,|z_{01}|^2),
\]
where $d_i$ are the diagonal entries and $z_e$ the cyclic off-diagonal entries.
The full relative-phase quotient also retains
$t=z_{12}z_{20}z_{01}$, with $|t|^2=\prod_e|z_e|^2$; it is generically
seven-dimensional. Thus the six displayed coordinates are a strain shadow
that discards the triangle phase, not the entire quotient
\cite{FrozenStrain}.

If $J=D\Phi$ and $G_Q$ is the reduced Green of a positive source Hessian
on its harmonic complement, minimizing source cost at fixed strain gives
$H_{\rm eff}=(JG_QJ^*)^{-1}$ where the reduced covariance is invertible.
For example, observing
only TT strain means using $J_{\rm TT}=P_{\rm TT}J$ in this formula.
Inverting the full covariance and then projecting generally gives a different
operator. Under a supplied same-source $S_3\times C_4$ action preserving
domains and $G_Q$, fixing or transporting the background, and intertwining
$J$ with $J_{DE}$, the two equivalent TT copies acquire the scalar Reynolds
shape. The symmetry lift and overall action scale are not selected by this
conditional implication.

The later support-stratified theorem covers the global cone
$\mathbb R^3\times[0,\infty)^3$ \cite{FrozenSupport}. With support $S$,
$\operatorname{rank}D\Phi=3+|S|$. A zero edge first appears as
$|\delta z_e|^2$, so its source quadratic cost is degree one in the normal
strain coordinate. Tangential shorting remains valid, but no ordinary full
six-dimensional boundary Hessian exists; a pseudoinverse does not invent
those missing normal directions. Rank-six surjectivity or the required local
$C_4$ derivative forces the regular all-nonzero-edge stratum. Global
regularity therefore is not an additional independent physical input once
that contract holds. The selected endpoint, reduced Green and physical
symmetry lift remain separate obligations. $J_{DE}$ is a tangent operation,
not an automorphism of the entire orthant.

\subsection{Finite transfer is not a physical strain action}
\label{sec:consumer-transfer}
Cohesive Closure Repair owns the exact symmetric Weyl transfer
\cite{CohesiveCurrent,FrozenMThree,FrozenMFour,FrozenHigher,FrozenNaturality}.
Its 144-dimensional DGA retracts to the old 36-dimensional complex together
with a twelve-dimensional higher-jet harmonic ideal. The resulting
48-dimensional target is not a minimal model: its differential is nonzero,
and its nonassociative binary product is accompanied by nonzero higher
operations. The complete $m_3$ and $m_4$ results are refined by a proved
nonzero family at every arity, not by a complete table from $m_5$ onward.
Contraction-preserving source maps transport all arities by the tree formula;
finite probes alone would not prove that conclusion. None of these
operations is a physical interaction vertex without the selected action and
continuum comparison. Fixed Points~I remains the owner of the analytic
existence and contraction gates: applying them here still requires a selected
operator, invariant complete domain, and the relevant gap, derivative and
resolvent bounds \cite{FixedPointsCurrent}. The finite transfer supplies none
of those physical inputs by itself.

\subsection{Cost pullback does not imply operator transport}
\label{sec:cost-versus-intertwiner}
The preceding conditions are stronger than agreement of a restricted
quadratic cost. For an elementary example take
\[
 Ux=(x,0),\qquad
 H_c=\begin{pmatrix}1&1\\1&2\end{pmatrix},\qquad H_f=(1).
\]
Here $H_c$ is positive definite and $U^*H_cU=H_f$, so the cost pulls back
exactly. Nevertheless $H_cUx=(x,x)$ whereas $UH_fx=(x,0)$: the selected
range is not invariant. Even exact intertwining on one invariant subspace
does not exclude additional zero modes in an unrepresented orthogonal
subspace. These are bounded linear-algebra illustrations of the interface,
not new physical models.

The Proto-Spinor paper, version~8, explains the source-owned overlap-polar
construction and its domain, scale, leakage and complement conditions
\cite{ProtoSpinorCurrent}. Those conditions are what allow finite Hessian
information to be transported; agreement of dimensions or positive repair
cost is insufficient. A positive repair functional also must not be
substituted for a physical signed action.

\subsection{What the exact internal TT support supplies}
\label{sec:gr-tt-support}
The frozen GR/QG source addresses a different but compatible question: which
internal directions can couple to the physical transverse-traceless (TT)
quotient on the selected exact branch? That quotient is the real two-plane
of plus and cross polarizations. A rotation of the transverse frame acts
with twice its angle. In the central-circle regular carrier
$\mathbb C[\mathbb Z_{64}]$, the corresponding real plane is spanned by the
cosine and sine modes $c_2,s_2$, equivalently the conjugate characters
$k=2,62$. Selecting the exact $d_*$ branch fixes the internal carrier; it
does not select an SI length or an arbitrary direction inside the strain
flag \cite{FrozenTT}.

The source's proof chain first accepts the selected core co-shape as
$B_0^*P_{\rm TT}=U_{\rm TT}C$, where $C$ is an invertible change of TT
basis and inner-product normalization. Its canonical value is $I_2$, not a
physical fit parameter. The finite support calculation supplies rank two,
zero leakage and same-angle central-shift intertwining. The selected
proper-time factorization then transports this core support to the dressed
map. With these premises the exact conclusion is
\[
 \Pi_{\rm exact64}B^*P_{\rm TT}=B^*P_{\rm TT},\qquad
 \operatorname{supp}(J_{\rm TT})
 =|d_*\rangle\otimes\operatorname{span}_{\mathbb R}\{c_2,s_2\}.
\]
The branch tower supplies $\lambda_{{\rm GR},{\rm TT}}=15$ in normalized
internal units. This is the accepted scoped source theorem, not an
independent verification inferred from boolean certificate fields.

There are two important type distinctions. First, this helicity-two support
is not the whole six-dimensional closure-strain space or its selected
three-dimensional shear sector. A relation requires the physical source map,
not matching ranks. Second, the value $15$ is not automatically the $\kappa$
in the finite q79 repair Hessian, a graviton mass, Newton's constant, or a
Planck scale. The full physical stress-energy response and dimensionful
normalization remain separate. Globally the external helicity bundle over
momentum directions must also remain distinct from the internal flat shared
line. The internal support result is retained exactly; these qualifications
identify its consumers rather than downgrading it.

\section{Status ledger}

\begin{center}
\small
\renewcommand{\arraystretch}{1.16}
\begin{tabularx}{\textwidth}{@{}
 >{\raggedright\arraybackslash}p{0.31\textwidth}
 >{\raggedright\arraybackslash}p{0.20\textwidth}
 >{\raggedright\arraybackslash}X@{}}
\hline
Object & Status & Meaning\\
\hline
$9=3+6$ comparison split & proved & local orientation/strain decomposition\\
$6=1+2+3$ strain split & proved with flag & exact orthogonal projectors\\
Iwasawa $SO(3)AN$ dimensions & proved & local/group normal form\\
q79 trace-split rank carrier & selected theorem & global rank $1+2+3$ carrier\\
q79 shared line and finite Hessian & proved at finite-symbol tier & connection,
holonomy, projector, and Hessian square\\
Internal TT co-shape support & exact on selected branch & helicity-two support;
internal eigenvalue $15$, not SI normalization\\
Finite SM algebra and profile operator & closed at declared profile tier & embedded renormalized-SM equivalence\\
Unique Higgs from Hessian & not a theorem & requires representation and alignment source\\
Three families from three lanes & not a theorem & requires index/monodromy source\\
Local strain--q79 identification & conditional/open & nonlinear shadow exact;
physical endpoint, Green and symmetry lift required\\
Strict no-knob value selection & stronger upgrade & not claimed here\\
\hline
\end{tabularx}
\end{center}

\section{Conclusion}

Closure-strain geometry supplies a useful and exact six-dimensional local
normal form.  Its real achievement is the explicit $1+2+3$ decomposition and
its compatibility target with the selected q79 carrier. The target is now
stronger than a rank match: its common flat differential line and normalized
finite Hessian square are exact. Standard Model
organization becomes credible only when the finite representation and source
packets are cited at their actual tier.  The paper therefore replaces broad
inevitability claims with one concrete bridge theorem capable of promoting the
local geometry into the already executed numerical branch.

% BEGIN MTT MANAGED COMPUTATIONAL EVIDENCE
\section*{Computational Evidence and Reproducibility}
The following earlier capsule is retained as provenance, not as a current
count of open obligations. New contextual imports use the immutable frozen
citations in the bibliography and the boundaries stated in this revision.
The local closure-strain decomposition is proved in the paper. The selected q=79, finite-matrix, anomaly, HYM, and internal TT rows provide concrete realizations or compatibility checks; they do not turn the conditional matter interpretation into a source theorem. The remaining Standard Model packets are corpus-state context.

The referenced rows are frozen to the curated results repository state identified below. Hashes are grouped in eight-character blocks for line breaking.
\begin{quote}\small
\textbf{Repository:} \url{https://github.com/PeterNero/mtt-results-repro}\\
\textbf{Commit:} \texttt{31247ebb 5c22f3fb b5443024 365433c6 ee0bff4a}\\
\textbf{Manifest:} \path{release/result_manifest.json}\\
\textbf{Manifest SHA-256:}\\
\texttt{fb399689 60b00584 631dbf53 1a708e18}\\
\texttt{ef928d6b 6d935119 c185d7f6 32b1e7cd}
\end{quote}

Tier labels are quoted verbatim from that manifest. A row used directly supports only the specific computational statement identified above; a corpus-state cross-check does not prove this paper's local theorems; and an open row is evidence of an unresolved obligation, never of closure.

\paragraph{Rows used directly in this paper.}
\begin{itemize}
\setlength{\itemsep}{0.25em}
\item \path{A22/e6_qpsi_qcd_anomaly} (\emph{derived exact}).\par E6 Qpsi matter/exotic QCD anomaly cancellation audit.
\item \path{A13/gr_tt_support} (\emph{derived exact}).\par Exact-branch internal TT support certificate; physical normalization remains open.
\item \path{A19/hym_wiener_contraction} (\emph{numeric certified}).\par Weighted-theta Fourier-tail and Wiener contraction certificate.
\item \path{A11/q79_exact_audit} (\emph{derived exact}).\par Executable q=79 exact-branch audit.
\item \path{A11/q79_exact_theorem} (\emph{derived exact}).\par CRT q=79 theorem on the selected exact branch.
\item \path{A01/qutrit_weyl_27_matrix} (\emph{derived exact}).\par Sparse 27x27 qutrit-Weyl left-action realization.
\end{itemize}

\paragraph{Corpus-state cross-checks.}
\begin{itemize}
\setlength{\itemsep}{0.25em}
\item \path{A01/charged_yukawa_higgs_profile} (\emph{profile replay}).\par Versioned Yu, Yd, Ye and lambda\_H profile packet.
\item \path{A14/ckm_prediction_profile} (\emph{numeric certified}).\par Three selected CKM profile rows and uncertainty comparison.
\item \path{A01/current_global_lock} (\emph{profile replay}).\par Current non-looping global status and source-certificate map.
\item \path{A01/direct_k_higgs_row} (\emph{derived exact}).\par Promoted direct K\_threshold.Omega\_H.lambda row.
\item \path{A04/final_12_of_12_audit} (\emph{profile replay}).\par Twelve-obligation embedded renormalized-SM equivalence audit.
\item \path{A40/neutral_two_primitive_profile} (\emph{profile replay}).\par Measured two-splitting neutrino profile, masses and Dirac Yukawa rows.
\item \path{A02/precision_15_source_transport} (\emph{profile replay}).\par Fifteen measured source coordinates, Jacobian and covariance transport.
\item \path{A02/precision_8x8_workspace} (\emph{profile replay}).\par Eight-coordinate SMDR output with positive-definite 8x8 covariance.
\item \path{A01/strict_pew_row} (\emph{derived exact}).\par Promoted P\_EW source row at the declared one-shared-primitive standard.
\end{itemize}

\paragraph{Open boundary (not evidence of closure).}
\begin{itemize}
\setlength{\itemsep}{0.25em}
\item \path{A05/strict_upgrade_ledger} (\emph{open}).\par Current 2/9 strict no-knob upgrade ledger.
\end{itemize}

No imported row changes theorem ownership or promotes a neighboring claim: all local statements retain their stated hypotheses, domains, and limitations.
% END MTT MANAGED COMPUTATIONAL EVIDENCE

\begin{thebibliography}{99}
\bibitem{QMGeometry} P.~Nero, \emph{Modal Triplet Theory: From MTT to Quantum Mechanics}, projective HYM, \v{C}ech compiler, spectral-strain and support-stratification sections, current owner manuscript, 2026.
\bibitem{CohesiveCurrent} P.~Nero, \emph{Cohesive Closure Repair and Its Hodge, Kernel, and Projection Shadows}, version~14, response-retract and all-arity source-map sections, 2026. The structural results are retained in the additive successor.
\bibitem{FixedPointsCurrent} P.~Nero, \emph{Fixed Points I: Projected Fixed Points, Equilibria, and Noncompact Domains in Modal Triplet Theory}, current analytic owner manuscript, 2026.
\bibitem{FrozenProjective} P.~Nero, \frozenresult{q79_projective_hym_naturality}{Covariant projective-module HYM naturality}, frozen result, 2026.
\bibitem{FrozenCech} P.~Nero, \frozenresult{q79_cech_projector_compiler}{Explicit curved connection compiler}, frozen result, 2026.
\bibitem{FrozenStrain} P.~Nero, \frozenresult{q79_spectral_strain_hessian}{Intrinsic spectral-strain quotient and shorted Hessian}, frozen result, 2026.
\bibitem{FrozenSupport} P.~Nero, \frozenresult{q79_support_stratified_strain}{Global support-stratified strain and boundary-Hessian exclusion}, frozen result, 2026.
\bibitem{FrozenMThree} P.~Nero, \frozenresult{causal_base_q79_symmetric_response_retraction_transferred_m3_packet}{Symmetric response retract and transferred ternary operation}, frozen result, 2026.
\bibitem{FrozenMFour} P.~Nero, \frozenresult{causal_base_q79_symmetric_response_transferred_m4_packet}{Transferred quaternary operation and arity-four coherence}, frozen result, 2026.
\bibitem{FrozenHigher} P.~Nero, \frozenresult{causal_base_q79_higher_transfer_jet_filtration_and_m5_feasibility_packet}{Higher-jet support and all-arity nontruncation}, frozen result, 2026.
\bibitem{FrozenNaturality} P.~Nero, \frozenresult{causal_base_q79_all_arity_source_promotion_packet}{All-arity contraction-morphism naturality}, frozen result, 2026.
\bibitem{Helgason}
S.~Helgason, \emph{Differential Geometry, Lie Groups, and Symmetric Spaces},
Academic Press, 1978.

\bibitem{Connes}
A.~Connes, \emph{Noncommutative Geometry}, Academic Press, 1994.

\bibitem{MTTSM}
P.~Nero, \emph{MTT Current True SM Closure Consolidated Ledger}, internal
theorem and verification packet, 2026.

\bibitem{SharedLine}
P.~Nero, \emph{q79 Universal Shared Differential Line and Finite-Operator
Intertwiner}, executable theorem packet, 2026.
\bibitem{FrozenTT}
P.~Nero, \href{https://github.com/PeterNero/mtt-results-repro/blob/f141a20ea23c5c3ff19cc2161c0e226e29ade8a7/release/results/gr_tt_support/artifact.json}{GR TT Support Final Theorem}, frozen exact-branch result, 2026.
The version~8 contextual import uses this immutable snapshot; the older
managed evidence snapshot above remains an earlier reference.
\bibitem{ProtoSpinorCurrent}
P.~Nero, \emph{The Proto-Spinor: Conditional Spinorial Closure and the q79
Interface}, version~8, 2026, sections on finite operators and the continuum
construction. Owner exposition of the frozen FSB.03a--e results.
\end{thebibliography}



\end{document}
